% Generated draft: requires mathematical and source review.
\documentclass[11pt]{article}
\usepackage[margin=25mm]{geometry}
\usepackage{fontspec}
\setmainfont{Latin Modern Roman}
\usepackage{amsmath,amssymb,mathtools,unicode-math}
\setmathfont{Latin Modern Math}
\usepackage{enumitem,needspace,xurl,hyperref}
\setlength{\parindent}{0pt}
\setlength{\parskip}{5pt}
\setlength{\emergencystretch}{3em}
\newcommand{\N}{\mathbb{N}}
\newcommand{\Z}{\mathbb{Z}}
\newcommand{\Q}{\mathbb{Q}}
\newcommand{\R}{\mathbb{R}}
\newcommand{\C}{\mathbb{C}}
\newcommand{\F}{\mathbb{F}}
\newcommand{\A}{\mathbb{A}}
\newcommand{\PP}{\mathbb{P}}
\newcommand{\E}{\mathbb{E}}
\newcommand{\eps}{\varepsilon}
\newcommand{\dd}{\,\mathrm{d}}
\begin{document}
\section*{UnsolvedMath ID: 30004399}
\textbf{Source problem code:} catalogue entry
\subsection{Definitions and mathematical statement}
Let \(G\curvearrowright X\) denote an action of a group \(G\) on a topological space \(X\) by homeomorphisms, and let \(D\subseteq X\) denote a chosen fundamental domain in the sense intended by the source. Necessarily \(G D=X\), where \(G D=\bigcup_{g\in G}gD\). For terminology, a strict fundamental domain meets each orbit in exactly one point:
\[
\lvert D\cap Gx\rvert=1\qquad\text{for every }x\in X,
\qquad Gx=\{gx:g\in G\}.
\]
Under this convention, a fundamental domain is non-strict if some orbit meets it in more than one point. Noncompactness refers to the subspace topology on \(D\). These conventions supply notation only; the record does not specify any further topological or combinatorial requirements on a fundamental domain, including its boundary identifications.

\textbf{Question (source formulation).} Under what hypotheses on the action \(G\curvearrowright X\) and its chosen fundamental domain \(D\) can the construction of the Bestvina complex be generalized when \(D\) is noncompact, non-strict, or both?

\textbf{Unresolved specification.} The supplied record neither defines the Bestvina complex nor states the construction to be extended or the properties its extension must retain. It also does not specify the class of groups, spaces, or actions under consideration. Consequently, the question above preserves the supplied research target, but a precise existence predicate or characterization cannot be extracted from this record alone. In particular, discreteness of \(G\), properness of the action, contractibility of \(X\), cellular structure, and restrictions on stabilizers are not supplied hypotheses and are not imposed here. The question concerns a chosen fundamental domain; the record does not specify whether the desired criterion should instead quantify over all, or over some, fundamental domains for a fixed action.
\subsection{Short English statement}
The problem asks when the Bestvina complex construction extends to actions whose chosen fundamental domain is noncompact, non-strict, or both. The intended construction, required output properties, and hypotheses on the action are unspecified, so a precise characterization requires further source clarification. This is a statement-only draft requiring human mathematical review.
\subsection{Sources}
\begin{itemize}
\item Supplied local UnsolvedMath registry record: permanent ID \(30004399\), title \emph{Bestvina Complexes Beyond Compact Strict Fundamental Domains}, source problem code \texttt{catalogue entry}, in \emph{Oberwolfach Reports Open Problems}. The record gives proposed year 2020 and names no proposer.
\item Supplied report citation: \emph{Manifolds and Groups (2021)}, Oberwolfach Reports, DOI \texttt{10.4171/owr/2020/7}, pp.~417--463. \url{https://doi.org/10.4171/owr/2020/7}. The background attributes the problem collection to workshop open-problem discussions. The supplied report title year is retained alongside the DOI as given.
\item Supplied literature reference: Petrosyan--Prytula 2020. \url{https://doi.org/10.4171/GGD/581}. No article title, pinpoint citation, or construction details are supplied.
\item Literature assessment dated exactly \texttt{2026-08-22}: the supplied curation labels the problem open and reports that the checked literature treats noncompact or non-strict fundamental domains as further questions, with no general resolving construction found. This is supplied dated curation, not an independently verified literature assessment; its status and cited claims require independent verification.
\item The registry labels the statement \texttt{corrected\_verified} and describes editorial repairs preserving mathematical content. This supplied label does not independently verify the formulation or resolve its missing hypotheses.
\end{itemize}
\end{document}
